\section{How many vacancies could one referral fill?}\label{sec:remedy}
With tentative offers in hand, could the platform fill a vacancy by referring one unmatched applicant to it? We consider only affordable, feasible apartments or eligible courses already on the original report. Each referral must preserve the existing contracts together in a stable allocation.
\subsection{One coordinated referral after offers}\label{sec:two-referral}

Suppose the platform now observes an assignment $M_0$ stable on the initial lists. Join each unmatched applicant to each mutually acceptable program with a current vacancy. These edges are hidden: if one were displayed it would block $M_0$. Give each program its remaining capacity and let $r$ be the size of a maximum matching in this graph. With lists of $\min(k,d_i)$ programs, where $k\ge1$ and $d_i$ is feasible degree, a hidden option can always replace a displayed one. If an entitlement is protected outside the quota, the substitution is made in a discretionary slot. In both replays, every selected applicant had a full five-slot discretionary list.

\begin{theorem}\label{thm:two-referral}
Among extensions that keep every initial contract in a single stable assignment, the maximum number of additional recipients is $r$. One shortlist substitution per selected applicant attains it. Every stable outcome on the revised display gives each consenting applicant their named referral, even if others decline. Under applicant-proposing DA, every original recipient also receives a weakly preferred assignment.
\end{theorem}
\begin{proof}
The new contracts in any extension must form a residual matching. Conversely, choose any such matching $A$. For each consenting participant substitute their assigned program for their last discretionary displayed option. Write $W$ for the original assignment together with the accepted referrals. All new pairs in the display belong to $W$. On old edges the applicant weakly improved. An originally full program has the same holders. At an originally slack one, every applicant with an old edge weakly preferred their original assignment, by stability. Hence $W$ is stable on this display. DA weakly dominates $W$ for applicants. All stable selections have the same matched people \citep{roth1986}. If some stable selection omitted an accepted referral $e$, it would keep its recipient matched even after deleting that unused edge. On that same graph $W\setminus\{e\}$ is stable by the preceding argument, since this applicant now has a subset of their original list. That would leave the person unmatched in one stable outcome only, a contradiction.
\end{proof}
The maximum fixes the original contracts in a stable witness. Opening new options to incumbents would be a different intervention. The empirical calculation also requires deciding which hidden options a person would actually accept. We report both a top-five rank limit and all mutually acceptable vacancies. Work on facilitating matches, adding school seats and guaranteed school referrals provides related interventions \citep{facilitate,tworound,lepesure}.


\begin{table}[ht]\centering\small
\caption{Largest number of named referrals compatible with retaining all existing contracts in one stable allocation. NYC uses simulated Haiku offers. Chile uses rank-informed pooled DA offers. All counts use instances in their own row. A vacancy must be at an affordable, feasible listing or an eligible course already on the pupil's reported list.}\label{tab:ref}
\begin{tabular}{@{}lrrrr@{}}\toprule
& Vacant seats & Unplaced with an option & Any rank & Top five\\\midrule
NYC ordinary & 7,972 & 757 & 683 & 204\\
NYC enriched & 5,990 & 2,527 & 1,926 & 370\\
Chile 2022 & 601,565 & 128 & 127 & 0\\
Chile 2023 & 640,751 & 97 & 97 & 1\\\bottomrule
\end{tabular}
\end{table}

In ordinary NYC markets only 37 unplaced lower-income renter instances could use any vacancy. Of the maximum 683 additional placements, at least 22 and at most 25 could go to them. In enriched markets 1,601 unplaced lower-income renters could use a vacancy. Maximizing total referrals first, then minimizing or maximizing referrals to this group, gave 1,002--1,085 lower-income recipients among 1,926 placements. Both supply and the households sampled change across these designs.

Many residual school seats were in courses that no unmatched pupil had listed. With full take-up, at most 224 placements could be added across the two years. Just one was possible within the family's top five. Neither year had an eligible referral after the rank-blind allocation, so even after each policy's largest extension the rank-informed policy placed 761 fewer pupils. After uniform rental menus the corresponding maxima were 215 ordinary and 518 enriched placements.

We verified the school maxima with a second national flow implementation. We replaced each selected applicant's last discretionary option with their referral, reran DA, and checked each named assignment, the weak improvement for earlier recipients, the quota and continuity rights. A service would need to observe current offers, elicit acceptable alternatives and give each family a chance to decline. Table~\ref{tab:ref} reports the opportunities available under those conditions.
